\labsection{3}{Production Planning and Manufacturing in Odoo}{sec:lab03}

\begin{labproject}{Turn a production plan into products, a BOM, and a completed Manufacturing Order}
\textbf{Coverage.} Tutorial 3 + Chapter 5.\\
\textbf{Core system.} Odoo Products, Purchase, Inventory, and Manufacturing.\\
\textbf{Goal.} Connect production-planning concepts to the master data and transactions required to manufacture Road Bikes.

\begin{labmode}
Tutorial 3 asks for a production approach, forecast, and S\&OP plan. Odoo does not replace those planning calculations in this course. Instead, it demonstrates how an approved plan becomes executable master data and manufacturing transactions.
\end{labmode}

\medskip\noindent\textbf{Part A --- Choose the production approach.}

For your Tutorial 3 business venture, justify make-to-stock, make-to-order, or a more customised approach. For West End Bicycles, use the supplied guide's make-to-stock story: build finished bikes before the customer delivery.

\medskip\noindent\textbf{Part B --- Create finished products.}

Create Road Bike and, if required by the class, the additional sample finished products. For Road Bike:
\begin{itemize}
  \item Product Type = Goods;
  \item Track Inventory = enabled;
  \item Can be Sold = enabled;
  \item Sales Price = RM 1,000.
\end{itemize}

\medskip\noindent\textbf{Part C --- Create raw materials.}

Create Aluminium Frame, Rubber Tires, Brake Set, Gear Set, and Bicycle Chain as Goods with Track Inventory enabled and Can be Purchased enabled. Use the guide costs as reference values.

\medskip\noindent\textbf{Part D --- Build the Road Bike BOM.}

Use \pathmenu{Manufacturing \textrightarrow Products \textrightarrow Bills of Materials}. Create a BOM for 1 Road Bike with:
\begin{itemize}
  \item 1 Aluminium Frame;
  \item 2 Rubber Tires;
  \item 1 Brake Set;
  \item 1 Gear Set;
  \item 1 Bicycle Chain.
\end{itemize}

Before manufacturing, verify that components have actually been received into stock. The SCM lab explains purchasing in depth; for this lab, use stock created through the supplied purchase/receipt steps or the instructor's prepared database.

\medskip\noindent\textbf{Part E --- Manufacture 5 Road Bikes.}

Create a Manufacturing Order for quantity 5, assign the Production Planner as Responsible, confirm the order, check component availability, and complete production using the Odoo 19 flow described in the supplied guide. Mark the MO Done.

\medskip\noindent\textbf{Part F --- Verify the inventory transformation.}

Before completion, record component on-hand quantities. After completion, verify:
\begin{itemize}
  \item frame stock decreased by 5;
  \item tyre stock decreased by 10;
  \item each other component decreased by 5;
  \item Road Bike stock increased by 5.
\end{itemize}

This is the ERP evidence that the manufacturing event changed the shared inventory data.

\medskip\noindent\textbf{Part G --- Complete the Tutorial 3 planning work.}

Prepare the sales forecast and S\&OP table required by the tutorial, including the selected product mix and the 90--95\% utilisation requirement specified in that tutorial. Keep the planning units consistent.
\end{labproject}

\begin{labdeliverable}
Submit the Tutorial 3 presentation content plus screenshots of the Road Bike product, BOM, completed Manufacturing Order, and the before/after stock evidence. Your explanation must connect the BOM and MO to the production plan rather than treating them as unrelated Odoo forms.
\end{labdeliverable}
